%%%%%%%%%%%%%%%%%%%%%%%%%%%%%%%%%%%%%%%%%%%%%%%%%%%%%%%%%%%%%%%%%%%%%%%%%%%%%%%%
\documentclass[conference]{IEEEtran}

\usepackage[utf8]{inputenc}
\usepackage[T1]{fontenc}
\usepackage{graphicx}
\usepackage{booktabs}
\usepackage{amsmath}
\usepackage{hyperref}
\usepackage{subcaption}
\usepackage{multirow}

% tighten float spacing
\setlength{\textfloatsep}{4pt plus 1pt minus 1pt}
\setlength{\floatsep}{3pt plus 1pt minus 1pt}
\setlength{\intextsep}{3pt plus 1pt minus 1pt}
\setlength{\abovecaptionskip}{3pt}
\setlength{\belowcaptionskip}{2pt}

\title{\LARGE \bf
Unsupervised Learning and Dimensionality Reduction Report
}

\author{
\parbox{3in}{
\centering
Michael Ciniello \\
School of Computer Science \\
Georgia Institute of Technology \\
Atlanta, GA, USA \\
{\tt\small mciniello3@gatech.edu}
}
}

\begin{document}

\maketitle
\thispagestyle{empty}
\pagestyle{empty}


%%%%%%%%%%%%%%%%%%%%%%%%%%%%%%%%%%%%%%%%%%%%%%%%%%%%%%%%%%%%%%%%%%%%%%%%%%%%%%%%
\section{Introduction and Data}

In the supervised reports we trained classifiers with explicit labels. This
report removes labels from the learning objective and asks a different
question: what structure can unsupervised algorithms discover in the same
feature spaces, and does that structure help or hinder the downstream
classification task?

We reuse the same two datasets and preprocessing pipeline from the SL and
OL Reports; no material changes were made. Crucially, this includes the
same feature engineering and \texttt{StandardScaler} standardization
(zero mean, unit variance) fitted on the training set only. Retaining
the identical preprocessing ensures that the neural network experiments
in Steps~4--5 are directly comparable to SL/OL baselines, while
standardization is essential for distance-based algorithms (K-Means,
GMM) whose objectives are sensitive to feature scale.

\textbf{Adult Income} (binary classification). The UCI Adult
dataset~\cite{kohavi1996} contains 45{,}222 rows with 6 numeric and 8
categorical features. The target encodes whether annual income exceeds
\$50K: class~0 ($\leq$50K, 75.2\%) vs.\ class~1 ($>$50K, 24.8\%).
Preprocessing includes median imputation, \texttt{StandardScaler},
one-hot encoding with \texttt{max\_categories=20}, and feature engineering
(binary capital-gain/loss flags, \texttt{log1p} transforms for
zero-inflated distributions), yielding 88 input dimensions. EDA revealed
that the highest numeric--target correlation is \texttt{education-num}
($|r|=0.333$). The resulting feature space is sparse and high-dimensional
due to one-hot encoding, which influences both clustering and
dimensionality reduction behavior.

\textbf{Wine Quality} (multiclass classification). The combined red/white
Wine Quality dataset~\cite{cortez2009} contains 6{,}497 rows with 11
continuous chemical features and a binary type indicator (red/white, roughly
1:3). The target \texttt{quality} takes 8 ordinal values but is heavily
concentrated in classes 3--6 ($>$93\% of samples). EDA showed that adjacent
quality levels share nearly identical feature distributions, that
\texttt{total\_sulfur\_dioxide} has the strongest target correlation
($|r|=0.556$), and that the red/white type distinction induces marked
chemistry shifts (e.g., sulfur dioxide, residual sugar). Preprocessing
applies \texttt{StandardScaler} only, giving 12 continuous input
dimensions --- a compact, dense feature space that is more naturally
suited to distance-based clustering than Adult's sparse one-hot
representation.

%%%%%%%%%%%%%%%%%%%%%%%%%%%%%%%%%%%%%%%%%%%%%%%%%%%%%%%%%%%%%%%%%%%%%%%%%%%%%%%%
\section{Metrics}

Without a labeled objective to optimize, we evaluate from two angles:
label-free metrics assess geometric cluster quality on their own terms,
while label-alignment metrics measure post-hoc agreement with known
classes to gauge whether discovered structure is semantically
meaningful. For the neural network experiments we retain the supervised
metrics from prior reports.

\textbf{Clustering (label-free).} \emph{Silhouette score} measures how
similar each point is to its own cluster versus the nearest neighboring
cluster, ranging from $-1$ to $+1$ (higher is better); we subsample to
5{,}000 points on Adult for computational feasibility.
\emph{Calinski-Harabasz} (CH) is the ratio of between-cluster to
within-cluster variance, scaled by degrees of freedom (higher is
better). \emph{Davies-Bouldin} (DB) averages, over all clusters, the
worst-case ratio of intra-cluster scatter to inter-centroid distance
(lower is better). \emph{BIC} (Bayesian Information Criterion)
penalizes log-likelihood by a term proportional to the number of
free parameters and $\log n$, favoring parsimonious models (lower is
better).

\textbf{Clustering (label alignment).} \emph{Adjusted Rand Index}
(ARI) measures pairwise agreement between cluster assignments and true
labels, adjusted for chance (0 = random, 1 = perfect).
\emph{Normalized Mutual Information} (NMI) quantifies how much knowing
the cluster assignment reduces uncertainty about the true label (and
vice versa), normalized to $[0,1]$. Both
are used strictly post-hoc --- labels are never used for model
selection.

\textbf{Neural network (Steps 4--5, Adult only).} Accuracy and binary
F1, consistent with prior reports. All NN experiments report mean
$\pm$ std across five seeds.

\textbf{Metric--algorithm alignment.} We select $k$ by silhouette for
K-Means and by BIC for GMM, because each criterion matches its
algorithm's assumptions. K-Means partitions space into spherical
Voronoi cells of roughly equal radius, so silhouette and CH --- which
reward compact, convex, similarly-sized clusters --- are natural
evaluation metrics. GMM fits full covariance matrices that produce
ellipsoidal, variably-sized components; these valid solutions can
score poorly on silhouette and CH simply because they overlap in
Euclidean space, whereas BIC evaluates the model's likelihood directly
and accounts for its additional complexity. DB complements both by
focusing on pairwise cluster confusion, penalizing clusters that are
simultaneously loose and close together regardless of shape. Reporting
all three label-free metrics alongside BIC gives a more complete
picture across both algorithms.

%%%%%%%%%%%%%%%%%%%%%%%%%%%%%%%%%%%%%%%%%%%%%%%%%%%%%%%%%%%%%%%%%%%%%%%%%%%%%%%%
\section{Hypotheses}

\textbf{H1 (Cluster structure reflects geometry, not labels).}
Unsupervised cluster assignments will not align well with classification
labels on either dataset --- we predict ARI $< 0.2$, well below the
moderate-agreement range. Adult's income threshold is a
socioeconomic gradient that cuts across demographic subgroups rather
than forming a compact geometric boundary in feature space. Wine's
quality scores are ordinal with heavy class overlap --- adjacent
levels share nearly identical chemical profiles --- so clusters will
capture coarser structural variation (e.g., the red/white type
distinction) rather than fine-grained quality
differences.

\textbf{H2 (PCA preserves the most useful signal for downstream tasks).}
Among PCA, ICA, and RP at matched dimensionality, PCA will yield the
highest downstream neural network F1 because it maximizes variance
preservation~\cite{jolliffe2002}. In general this need not align with
class separation, but our EDA suggests it holds here: the features
most correlated with the target are also among the highest-variance
features after scaling, and the positive-class rate varies
monotonically with these predictors, so preserving variance should
preserve discriminative signal. ICA seeks axes of maximal statistical
independence (non-Gaussianity)~\cite{hyvarinen2000}; the resulting
components may isolate interesting latent sources (e.g., zero-inflated
capital features) but need not align with the decision boundary,
so we expect higher sensitivity to weight initialization and a
slightly larger F1 drop. RP preserves pairwise distances with high
probability per the Johnson-Lindenstrauss
lemma~\cite{dasgupta2003}, but random axes discard variance
indiscriminately, yielding the highest reconstruction error and a
meaningful F1 drop.

\textbf{H3 (Continuous cluster features outperform discrete ones).}
Adding cluster-derived features (centroid distances, GMM posteriors)
to the raw feature set will improve F1 more than adding one-hot
cluster labels, because continuous features encode graded information
--- how far a sample is from each centroid, or how likely it belongs
to each component --- that the NN can exploit directly, whereas
one-hot labels reduce all within-cluster variation to a single
category. Furthermore, cluster features used \emph{alone} (without
raw features) will capture meaningful but incomplete signal,
demonstrating that unsupervised decomposition recovers a nontrivial
fraction of the classification task even without access to the
original feature space.

%%%%%%%%%%%%%%%%%%%%%%%%%%%%%%%%%%%%%%%%%%%%%%%%%%%%%%%%%%%%%%%%%%%%%%%%%%%%%%%%
\section{Methodology}

\textbf{Clustering algorithms.} K-Means (\texttt{sklearn.cluster.KMeans},
$n\_init=10$) and Expectation Maximization via Gaussian Mixture Models
(\texttt{sklearn.mixture.GaussianMixture}, full covariance,
$n\_init=5$, $\text{reg\_covar}=10^{-4}$). K-Means uses Euclidean
distance and assigns each point to the nearest centroid; this is
well-suited to both datasets after standardization, since all features
are on a common scale. GMM generalizes K-Means by fitting full
covariance matrices~\cite{reynolds2009}, allowing ellipsoidal clusters and soft
(probabilistic) assignments. However, fitting full covariance on
Adult's 88-dimensional space requires estimating $O(d^2)$ parameters
per component, and the sparse one-hot features create columns that
are mostly zero with only a few distinct values. When a component
contains too few samples with nonzero values in a given feature, the
corresponding covariance entries become near-zero, driving the
covariance matrix toward singularity. Since the E-step requires
inverting each component's covariance matrix to evaluate the
multivariate Gaussian density, a singular $\Sigma_k$ causes the
computation to fail. We address this by raising
\texttt{reg\_covar} from the default $10^{-6}$ to $10^{-4}$, which
adds $\lambda I$ to each covariance matrix, shifting near-zero
eigenvalues to at least $\lambda$ and making the matrix safely
invertible --- analogous to ridge regularization in linear regression.
The value $10^{-4}$ is large enough to stabilize Adult's sparse
features while remaining small enough to preserve GMM's ability to
fit ellipsoidal components.
Wine does not require this adjustment because its 12 continuous
features have smooth, well-spread distributions that keep all
covariance matrices well-conditioned. We sweep $k \in [2,20]$,
selecting the best $k$ by silhouette score (K-Means) or BIC (GMM);
labels are never used for model selection but are referenced post-hoc
for alignment analysis.

\textbf{Dimensionality reduction.} PCA selects the number of components
at the 95\% cumulative-variance threshold, a widely used default that
balances dimensionality reduction against information
loss~\cite{jolliffe2002}. Neither ICA nor RP has an analogous variance-based selection
criterion --- ICA maximizes statistical independence, and RP relies on
probabilistic distance-preservation guarantees --- so we set both to
the same $n$ as PCA. This controls for dimensionality and enables
apples-to-apples comparison across methods in Steps~3--4; any
performance differences are attributable to the projection itself,
not the number of output features.
ICA uses \texttt{sklearn.decomposition.FastICA} with unit-variance
whitening. RP uses scikit-learn's \texttt{GaussianRandomProjection}
with reconstruction error computed via the pseudo-inverse of the
projection matrix and averaged over 10 random seeds to account for
the stochasticity of the random matrix.

\textbf{Neural network (Steps 4--5, Adult only).} The same architecture
and hyperparameters as the best OL model: \texttt{TabularMLP}
(256--128--64, ReLU, no dropout), Adam ($\alpha=0.001$,
WD$=10^{-3}$), batch size 256, up to 200 epochs with
early stopping on validation F1 (patience 25).

\textbf{Data splits.} We reuse the OL stratified 80 / 20 train/test split
(\texttt{random\_state=42}), with the training set further split 85/15
into train and validation subsets (Adult: 30{,}750 / 5{,}427 / 9{,}045;
Wine: 4{,}417 / 780 / 1{,}300). All unsupervised models are fit on training
data only; test sets are transformed but never used for fitting.

\textbf{Seeds and hardware.} Seeds [42, 123, 456, 789, 1024]; Apple M2
Pro, 16\,GB RAM, CPU backend (MPS disabled for reproducibility).

%%%%%%%%%%%%%%%%%%%%%%%%%%%%%%%%%%%%%%%%%%%%%%%%%%%%%%%%%%%%%%%%%%%%%%%%%%%%%%%%
\section{Step 1: Clustering on Raw Data}

\begin{table}[t]
\centering
\caption{Step~1 summary. K-Means selects $k{=}2$ on both datasets;
GMM selects much higher $k$ via BIC. Label alignment (ARI, NMI) is
weak across the board.}
\label{tab:step1}
\resizebox{\columnwidth}{!}{%
\begin{tabular}{llrcccccc}
\toprule
\textbf{Data} & \textbf{Algo} & $k$ & \textbf{Silh.} & \textbf{CH} & \textbf{DB} & \textbf{BIC} & \textbf{ARI} & \textbf{NMI} \\
\midrule
\multirow{2}{*}{Adult}
  & K-Means &  2 & .440 & 5446 & 0.91 & --- & .064 & .023 \\
  & GMM     & 19 & .057 & 1289 & 4.54 & $-$11.3M & .032 & .096 \\
\midrule
\multirow{2}{*}{Wine}
  & K-Means &  2 & .339 & 1688 & 1.34 & --- & .183 & .253 \\
  & GMM     & 12 & .022 &  346 & 3.47 & 57K & .098 & .205 \\
\bottomrule
\end{tabular}}
\end{table}

\begin{figure}[t]
\centering
\begin{subfigure}{0.48\columnwidth}
  \includegraphics[width=\linewidth]{results/step1/adult_kmeans_selection.png}
  \caption{Adult K-Means}
\end{subfigure}\hfill
\begin{subfigure}{0.48\columnwidth}
  \includegraphics[width=\linewidth]{results/step1/wine_kmeans_selection.png}
  \caption{Wine K-Means}
\end{subfigure}
\caption{K-Means selection curves. Silhouette peaks at $k{=}2$ for both
datasets; inertia shows a gradual elbow but no strong case for higher $k$.}
\label{fig:step1_selection}
\end{figure}

Table~\ref{tab:step1} summarizes the best cluster configurations
selected by unsupervised criteria only.

\textbf{K-Means selects $k{=}2$ for both datasets}
(Fig.~\ref{fig:step1_selection}). All three geometric metrics agree:
silhouette is highest, CH is highest, and DB is lowest at $k{=}2$
(Table~\ref{tab:step1}), unanimously favoring the simplest partition.
On Adult, silhouette drops sharply
from 0.44 ($k{=}2$) to 0.33 ($k{=}3$) to 0.13 ($k{=}4$), then
plateaus around 0.12--0.14 for all higher $k$. The steep initial drop
means the dominant binary split in the 88-dimensional one-hot space is
clean, but adding a third cluster immediately breaks that separation;
the subsequent plateau indicates no further well-separated structure
at any granularity. Inertia decays smoothly with no sharp elbow: each
additional cluster shaves off a similar amount of within-cluster
variance, with no single split being dramatically more informative
than the last. This is a signature of data whose variance is spread
across many dimensions rather than concentrated in a few --- the same
property PCA's eigenvalue spectrum reveals in Step~2, where Adult's
eigenvalues decay gradually with no dominant subspace. On Wine, the pattern is similar: silhouette peaks at
$k{=}2$ (0.34) and drops to $\sim$0.25 by $k{=}3$, confirming the
red/white type distinction is the dominant axis of chemical variation.
A small silhouette bump at $k{=}4$ hints at a secondary 2$\times$2
structure (possibly red/white crossed with a quality split), but it
does not exceed the $k{=}2$ peak. Wine's inertia curve is steeper
than Adult's, consistent with its more concentrated feature space (12
continuous dimensions vs.\ 88 sparse).

\textbf{What do the $k{=}2$ clusters actually capture?} To move
beyond selection curves, we profiled each cluster by computing
$|\Delta^{(j)}| = |\mu_1^{(j)} - \mu_0^{(j)}|$ for every feature~$j$,
where $\mu_c^{(j)}$ is the mean standardized value of feature~$j$ in
cluster~$c$. Because all features share the same standardized scale,
$|\Delta|$ directly quantifies how much a feature differs between the
two groups K-Means found; a larger value means the feature is a
stronger discriminator of the split.
Fig.~\ref{fig:step1_profile} shows the per-cluster feature means for
the top-10 features ranked by $|\Delta|$.

On Adult, the split is driven almost entirely by
\texttt{capital-loss} ($|\Delta|=4.73$ for its log-transformed variant): Cluster~1
(4.7\% of training data) contains all samples with nonzero capital
loss (\texttt{has\_capital\_loss}$=$1.0), while Cluster~0 contains
the remaining 95.3\%. This is a direct consequence of the
zero-inflated distribution identified in EDA --- roughly 96\% of
samples have exactly zero capital loss, creating a tight, separable
minority cluster that K-Means latches onto. Notably, this cluster is
51.4\% high-income vs.\ 23.5\% in the main cluster, explaining the
modest ARI: the capital-loss split partially correlates with income
(investment losses imply investment activity, which correlates with
wealth) but is not the income boundary itself. This motivates
applying dimensionality reduction before clustering (Step~3), where
PCA may dilute the outsized influence of individual zero-inflated
features and reveal broader structure.

On Wine, the top-ranked feature is \texttt{type} ($|\Delta|=2.30$).
Because \texttt{type} is a binary red/white indicator, we can verify
the split directly: Cluster~0 (3{,}328 samples) has a standardized
\texttt{type} mean of $+0.57$, nearly equal to the $+0.58$ expected if
the cluster were exclusively white, while Cluster~1 (1{,}089 samples)
has a mean of $-1.73$, matching the $-1.73$ expected if exclusively
red. The cluster sizes also mirror
the 3:1 white-to-red ratio in the dataset, confirming that K-Means
recovers the red/white distinction almost perfectly. The next six highest-$|\Delta|$ features are exactly the chemistry
features that EDA flagged as differing sharply between red and white
wines.
Because quality scores differ between wine types, the red/white
split incidentally carries information about the target variable
\texttt{quality}: 75\% of Cluster~0 (predominantly white) has
\texttt{quality} ratings of 3--4, while 83\% of Cluster~1
(predominantly red) has ratings of 5--6. This explains the moderate
NMI (0.253) with respect to the quality target, even though K-Means
never saw labels.

\begin{figure}[t]
\centering
\begin{subfigure}{\columnwidth}
  \includegraphics[width=\linewidth]{results/step1/adult_cluster_profile.png}
  \caption{Adult}
\end{subfigure}
\vspace{4pt}
\begin{subfigure}{\columnwidth}
  \includegraphics[width=\linewidth]{results/step1/wine_cluster_profile.png}
  \caption{Wine}
\end{subfigure}
\caption{Top 10 features by absolute centroid difference for K-Means
($k{=}2$). Adult's split is dominated by the zero-inflated
capital-loss feature; Wine's split captures the red/white type
distinction.}
\label{fig:step1_profile}
\end{figure}

\begin{figure}[t]
\centering
\begin{subfigure}{0.48\columnwidth}
  \includegraphics[width=\linewidth]{results/step1/adult_gmm_selection.png}
  \caption{Adult GMM}
\end{subfigure}\hfill
\begin{subfigure}{0.48\columnwidth}
  \includegraphics[width=\linewidth]{results/step1/wine_gmm_selection.png}
  \caption{Wine GMM}
\end{subfigure}
\caption{GMM selection curves. BIC improves monotonically before
plateauing at $k{=}19$ (Adult) and $k{=}12$ (Wine). Silhouette
remains low throughout, indicating overlapping components.}
\label{fig:step1_gmm_selection}
\end{figure}

\textbf{GMM selects much higher $k$ (Adult=19, Wine=12).} Whereas
K-Means collapsed to $k{=}2$ on both datasets, latching onto a single
dominant split, GMM's full-covariance components can model
sub-populations of varying shape and spread, so BIC continues to
reward additional components well beyond $k{=}2$.
As shown in Fig.~\ref{fig:step1_gmm_selection}, Adult's BIC decreases
steeply from $k{=}2$ to $k{=}9$, capturing the most important
distributional sub-populations, then continues to improve more slowly
through $k{=}19$ before plateauing. Wine's BIC shows a clearer elbow,
flattening around $k{=}8$--$12$; notably, AIC continues to decline
past $k{=}12$, but BIC's stronger complexity penalty halts there ---
illustrating how the choice of information criterion influences model
complexity. Raw BIC values differ by orders of magnitude between
datasets (millions for Adult, thousands for Wine) because BIC scales
with sample count and dimensionality; it is only meaningful as a
relative comparison across $k$ on the same dataset. Silhouette
remains low throughout both sweeps, but the two datasets behave
differently: Adult's silhouette is nearly flat and noisy
(0.03--0.10), reflecting how its 88 sparse dimensions make all
components overlap in Euclidean space, while Wine's spans a wider
range (0.02--0.34) with more visible structure, consistent with its
compact 12-dimensional feature space. In both cases, silhouette
tells a fundamentally different story than BIC: while BIC rewards
likelihood improvements from additional components, silhouette
reveals that those components overlap geometrically. CH and DB tell the same geometric story: GMM's selected $k$ yields
much lower CH and higher DB than K-Means' $k{=}2$
(Table~\ref{tab:step1}), confirming that GMM's components overlap
heavily by Euclidean measures. This divergence between geometric
metrics and BIC underscores why we use BIC for GMM selection ---
silhouette, CH, and DB all assume compact, convex clusters, which is
misaligned with GMM's probabilistic model.

\textbf{Label alignment is weak across the board} (ARI $< 0.2$
everywhere; Table~\ref{tab:step1},
Fig.~\ref{fig:step1_alignment}). This is expected: Adult income is a
socioeconomic threshold, not a natural cluster boundary in
demographic feature space, and Wine quality is ordinal with heavy
class overlap in physicochemical features.

On \emph{Wine}, both algorithms start with their highest alignment at
$k{=}2$ (K-Means ARI$=$0.183, NMI$=$0.253) and decline steadily as
$k$ increases --- additional clusters fragment the red/white split
that drives alignment, without uncovering new label-relevant
structure. Wine K-Means at $k{=}2$ achieves the highest alignment of
any configuration across both datasets, because the red/white split
correlates with quality distributions (as shown in the profiling
above).

On \emph{Adult}, the pattern differs. K-Means ARI spikes at $k{=}3$
($\approx$0.19, nearly at our 0.2 threshold) before dropping
steadily, suggesting three clusters briefly align with income but
fragment at higher $k$. GMM ARI at $k{=}2$ is
actually \emph{negative} ($-$0.05), meaning worse-than-random label
agreement --- its two components do not correspond to income at all.
However, GMM's NMI steadily rises with $k$ (reaching 0.096 at
$k{=}19$ vs.\ K-Means NMI$=$0.023 at $k{=}2$), suggesting that
additional components capture income-correlated subpopulations
without disrupting existing ones.

\begin{figure}[t]
\centering
\begin{subfigure}{\columnwidth}
  \includegraphics[width=\linewidth]{results/step1/adult_label_alignment.png}
  \caption{Adult}
\end{subfigure}
\vspace{4pt}
\begin{subfigure}{\columnwidth}
  \includegraphics[width=\linewidth]{results/step1/wine_label_alignment.png}
  \caption{Wine}
\end{subfigure}
\caption{Label alignment (ARI/NMI) vs.\ $k$ for K-Means and GMM.
Alignment is universally low, confirming that cluster structure does
not map onto classification labels.}
\label{fig:step1_alignment}
\end{figure}

%%%%%%%%%%%%%%%%%%%%%%%%%%%%%%%%%%%%%%%%%%%%%%%%%%%%%%%%%%%%%%%%%%%%%%%%%%%%%%%%
\section{Step 2: Dimensionality Reduction}

\begin{table}[t]
\centering
\caption{Step~2 summary. All three DR methods share matched output
dimensionality per dataset (set by PCA's 95\% variance threshold).}
\label{tab:step2}
\begin{tabular}{llrrr}
\toprule
\textbf{Data} & \textbf{Method} & \textbf{$d_{\text{in}}$} & \textbf{$d_{\text{out}}$} & \textbf{Key metric} \\
\midrule
\multirow{3}{*}{Adult}
  & PCA & 88 & 30 & 95.3\% var.\ retained \\
  & ICA & 88 & 30 & $\overline{|\kappa|}=11.2$ \\
  & RP  & 88 & 30 & recon.\ err.\ 0.816 \\
\midrule
\multirow{3}{*}{Wine}
  & PCA & 12 & 9 & 96.7\% var.\ retained \\
  & ICA & 12 & 9 & $\overline{|\kappa|}=14.0$ \\
  & RP  & 12 & 9 & recon.\ err.\ 0.478 \\
\bottomrule
\end{tabular}
\end{table}

Table~\ref{tab:step2} summarizes all three methods.

\textbf{PCA: strong compression for Adult, modest for Wine.} Adult
reduces from 88 to 30 components (66\% reduction) while retaining
95.3\% of variance. The first 5 PCs explain $\sim$60\%, reflecting
redundancy among one-hot columns. The scree plot
(Fig.~\ref{fig:step2_pca}) shows smooth decay with no sharp elbow
--- the signal is spread across many components. Wine reduces from 12
to 9 (25\%) at 96.7\% variance; with only continuous features, there
is less redundancy. Unlike Adult's smooth decay, Wine's scree plot
shows a clear elbow around PC3--4: the first three PCs alone explain
$\sim$67\%, suggesting variance is concentrated in a few meaningful
directions (PC1 at 31.8\% likely captures the red/white type
distinction; PC2--3 may reflect quality-correlated chemistry).

\begin{figure}[t]
\centering
\begin{subfigure}{\columnwidth}
  \includegraphics[width=\linewidth]{results/step2/adult_pca_variance.png}
  \caption{Adult PCA}
\end{subfigure}
\vspace{4pt}
\begin{subfigure}{\columnwidth}
  \includegraphics[width=\linewidth]{results/step2/wine_pca_variance.png}
  \caption{Wine PCA}
\end{subfigure}
\caption{PCA cumulative explained variance. The dashed line marks the
95\% threshold used for component selection.}
\label{fig:step2_pca}
\end{figure}

\textbf{ICA: high kurtosis confirms non-Gaussian structure.} Both
datasets show mean absolute kurtosis well above the Gaussian baseline
of 3 (Adult: 11.2, Wine: 14.0), confirming recoverable independent
sources. Adult ICA shows extreme kurtosis in some components
(max$=$169), driven by the zero-inflated capital-gain/loss features.
Wine's max kurtosis (78.9) is also high, likely reflecting the
bimodal red/white distribution.

\textbf{RP: significant information loss at matched dimensionality.}
At the same $n$, RP's relative reconstruction error (Frobenius norm
ratio) is 0.816 for Adult and 0.478 for Wine. This directly
illustrates PCA's optimality: PCA selects maximum-variance directions,
while RP picks random directions that inevitably miss signal.
Adult's nearly linear error curve (Fig.~\ref{fig:step2_rp}) reflects
its flat eigenvalue spectrum --- no small random subspace captures
the 88-dimensional spread. Notably, Wine's $\pm$1 std band is
visibly wider than Adult's, because with only 12$\to$9 dimensions each
random projection matrix has a larger impact on what signal is
preserved, whereas Adult's 88$\to$30 projection averages out across
more directions.

\begin{figure}[t]
\centering
\begin{subfigure}{0.48\columnwidth}
  \includegraphics[width=\linewidth]{results/step2/adult_rp_reconstruction.png}
  \caption{Adult RP}
\end{subfigure}\hfill
\begin{subfigure}{0.48\columnwidth}
  \includegraphics[width=\linewidth]{results/step2/wine_rp_reconstruction.png}
  \caption{Wine RP}
\end{subfigure}
\caption{RP reconstruction error vs.\ number of components.
Adult's error decreases nearly linearly (flat spectrum); Wine
drops faster (stronger dominant components).}
\label{fig:step2_rp}
\end{figure}

%%%%%%%%%%%%%%%%%%%%%%%%%%%%%%%%%%%%%%%%%%%%%%%%%%%%%%%%%%%%%%%%%%%%%%%%%%%%%%%%
\section{Step 3: Clustering on Reduced Data}

Each DR-transformed dataset was clustered with K-Means and GMM using
the same $k$-range and selection criteria as Step~1.
Table~\ref{tab:step3} compares all configurations against the raw-data
baselines.

\begin{table}[t]
\centering
\caption{Step~3: clustering on DR data vs.\ raw baselines. Bold marks
the best label-free score in each column per dataset.}
\label{tab:step3}
\begin{tabular}{llrccc}
\toprule
\textbf{Space} & \textbf{Algo} & $k$ & \textbf{Silh.} & \textbf{ARI} & \textbf{NMI} \\
\midrule
\multicolumn{6}{l}{\emph{Adult}} \\
RAW & K-Means &  2 & 0.440 & 0.064 & 0.023 \\
PCA & K-Means &  2 & 0.453 & 0.064 & 0.023 \\
ICA & K-Means & 11 & 0.133 & 0.016 & 0.044 \\
RP  & K-Means &  2 & \textbf{0.533} & \textbf{0.083} & \textbf{0.034} \\
RAW & GMM     & 19 & 0.057 & 0.032 & 0.096 \\
PCA & GMM     & 19 & 0.055 & 0.020 & 0.087 \\
ICA & GMM     & 18 & 0.099 & 0.019 & 0.050 \\
RP  & GMM     & 20 & 0.001 & 0.020 & 0.082 \\
\midrule
\multicolumn{6}{l}{\emph{Wine}} \\
RAW & K-Means &  2 & 0.339 & 0.183 & 0.253 \\
PCA & K-Means &  2 & \textbf{0.347} & \textbf{0.183} & 0.253 \\
ICA & K-Means &  2 & 0.176 & 0.181 & \textbf{0.258} \\
RP  & K-Means &  4 & 0.320 & 0.134 & 0.178 \\
RAW & GMM     & 12 & 0.022 & 0.098 & 0.205 \\
PCA & GMM     & 12 & 0.060 & 0.089 & 0.207 \\
ICA & GMM     & 15 & 0.055 & 0.087 & 0.203 \\
RP  & GMM     & 14 & 0.049 & 0.083 & 0.189 \\
\bottomrule
\end{tabular}
\end{table}

Fig.~\ref{fig:step3_heatmap} visualizes silhouette across all
combinations.

\textbf{PCA preserves raw-data structure.} On both datasets, PCA
K-Means selects the same $k{=}2$ with identical ARI and NMI to the
raw baseline, while silhouette improves slightly (Adult: 0.440$\to$0.453;
Wine: 0.339$\to$0.347). The matching metrics suggest that PCA's
variance-preserving projection retains the dominant cluster geometry,
with the removal of noise dimensions tightening clusters marginally.

\textbf{RP yields the best K-Means silhouette on Adult} (0.533)
and the highest ARI (0.083) --- a surprising result, since RP
discards the most information (reconstruction error 0.82). The likely
mechanism: in the original 88-dimensional space, silhouette is
computed across all dimensions, so the $\sim$85 sparse one-hot
columns that are irrelevant to the capital-loss split add noise to
distance calculations, diluting cluster separation. RP's random
mixing projects all 88 dimensions into 30 dense combinations; by the
Johnson-Lindenstrauss lemma, pairwise distances (and thus
between-cluster gaps) are approximately preserved, but the random
averaging smooths out sparse one-hot noise, reducing within-cluster
scatter relative to between-cluster separation. PCA does not get this
benefit because it faithfully preserves the highest-variance
directions, which include within-cluster variance from the one-hot
structure. Notably, this effect is K-Means-specific: on Adult, RP GMM
silhouette is essentially zero (0.001), as GMM selects $k{=}20$
overlapping components in just 30 dimensions --- consistent with the
BIC-vs-silhouette divergence discussed above.

\textbf{ICA fragments clusters.} On Adult, ICA K-Means selects
$k{=}11$, the highest of any method. ICA's independence-maximizing
rotation creates axes with high kurtosis (heavy tails, sharp spikes),
and these spiky structures look like small tight clusters to K-Means,
driving the algorithm to find many more partitions. Silhouette drops
from 0.440 to 0.133, yet NMI actually \emph{increases} (0.023
$\to$ 0.044): the 11 clusters capture more label-relevant
substructure than the coarse $k{=}2$ split, even though the
geometry is weaker. On Wine, ICA keeps $k{=}2$ with nearly identical label alignment
(ARI$=$0.181, NMI$=$0.258) but halved silhouette (0.339 $\to$ 0.176),
meaning ICA distorts geometry without disrupting the red/white
structure.

\textbf{GMM is largely insensitive to DR.} BIC selects similar $k$
across all feature spaces (Adult: 18--20, Wine: 12--15), suggesting
the data's intrinsic distributional complexity drives component count,
not the projection. Wine GMM NMI is remarkably stable (0.189--0.207
across all four spaces), meaning the quality-relevant information
captured by GMM components persists regardless of how the features
are transformed. On Adult, ICA is the one exception: GMM NMI drops
from $\sim$0.09 to 0.050, likely because ICA's kurtosis-maximizing
axes distort the subpopulation structure that GMM components
otherwise capture.

\begin{figure}[t]
\centering
\begin{subfigure}{0.48\columnwidth}
  \includegraphics[width=\linewidth]{results/step3/adult_silhouette_heatmap.png}
  \caption{Adult silhouette}
\end{subfigure}\hfill
\begin{subfigure}{0.48\columnwidth}
  \includegraphics[width=\linewidth]{results/step3/wine_silhouette_heatmap.png}
  \caption{Wine silhouette}
\end{subfigure}
\caption{Step~3: silhouette heatmaps across DR method $\times$
clustering algorithm. Higher is better. RP + K-Means dominates on
Adult; PCA + K-Means leads on Wine.}
\label{fig:step3_heatmap}
\end{figure}

%%%%%%%%%%%%%%%%%%%%%%%%%%%%%%%%%%%%%%%%%%%%%%%%%%%%%%%%%%%%%%%%%%%%%%%%%%%%%%%%
\section{Step 4: Neural Network on DR Data}

The same NN architecture as the baseline was trained on PCA-, ICA-,
and RP-transformed Adult data (30 dims, down from 88). Each
configuration was evaluated over 5 seeds.

\begin{table}[t]
\centering
\caption{Step~4 results (Adult). All DR methods incur a small F1
drop ($\sim$1~pp) for a 66\% dimensionality reduction.}
\label{tab:step4}
\begin{tabular}{lrccc}
\toprule
\textbf{Space} & \textbf{Dims} & \textbf{Accuracy} & \textbf{F1} & $\Delta$\textbf{F1} \\
\midrule
Baseline &  88 & .848$\pm$.003 & \textbf{.696}$\pm$.004 & --- \\
PCA      &  30 & .846$\pm$.002 & .687$\pm$.004 & $-$.009 \\
ICA      &  30 & .844$\pm$.002 & .684$\pm$.007 & $-$.011 \\
RP       &  30 & .843$\pm$.004 & .686$\pm$.003 & $-$.010 \\
\bottomrule
\end{tabular}
\end{table}

\textbf{All DR projections produce a small F1 drop ($\sim$0.9--1.1~pp)}
(Fig.~\ref{fig:step4_comparison}).
This is a modest cost for reducing dimensionality by 66\%. Accuracy
drops are even smaller ($<$0.5~pp), confirming that most predictive
information is preserved (Table~\ref{tab:step4}).

\textbf{PCA retains the most signal.} Its 95\%-variance threshold
preserves the linear structure that matters for classification. The F1
gap ($-$0.009) is within $\sim$2 standard deviations of baseline
variance, suggesting it could be bridged with minor tuning.

\textbf{RP is competitive with PCA despite high information loss.}
Despite reconstruction error $\sim$0.82, RP nearly matches PCA's F1.
The Johnson-Lindenstrauss lemma~\cite{dasgupta2003} offers an
explanation: RP approximately preserves pairwise distances, so the
projected data retains enough geometric structure for the NN to
recover similar decision boundaries. This echoes our Step~3 finding
that RP K-Means actually \emph{outperformed} PCA on Adult
(silhouette 0.533 vs.\ 0.453), though the mechanisms likely differ
(noise smoothing for clustering vs.\ general distance preservation
for classification).

\textbf{ICA has the largest drop and highest variance.} ICA optimizes for statistical independence, which is not guaranteed to
align with classification objectives --- the resulting axes may
emphasize distributional structure (e.g., heavy tails, spikes) that
does not discriminate income levels. The higher seed-to-seed variance ($\pm$0.007 vs.\ $\pm$0.004)
suggests that ICA's independence-optimized axes create a less smooth
loss landscape for the NN, making training more sensitive to weight
initialization.

\begin{figure}[t]
\centering
\includegraphics[width=0.95\columnwidth]{results/step4/adult_f1_comparison.png}
\caption{Step~4: F1 comparison (mean $\pm$ std across 5 seeds). All
DR methods fall slightly below baseline but the gaps are small.}
\label{fig:step4_comparison}
\end{figure}

\textbf{DR does not reduce total training time despite 66\% fewer
inputs} (Fig.~\ref{fig:step4_convergence}). All four methods reach a
similar val-F1 plateau ($\sim$0.69--0.70) within the first 10--15
epochs, confirming that DR features carry enough signal for rapid
learning. However, PCA trains significantly longer (mean 62 epochs
vs.\ 43 for baseline) because its smooth loss surface keeps yielding
marginal improvements that delay early stopping, without reaching a
higher final F1. ICA and RP converge at similar epoch counts to
baseline (43 and 44), consistent with their noisier loss landscapes
triggering patience sooner. Per-epoch wall time is marginally lower
for DR methods due to the smaller input layer, but the difference is
negligible on CPU --- the first linear layer (88$\to$256 vs.\
30$\to$256) is a small fraction of overall computation, so total
training time is dominated by epoch count rather than per-epoch cost.

For this dataset and architecture, DR is a viable trade-off: a
$\sim$1\% F1 sacrifice buys 66\% fewer input features with comparable
or faster convergence. The small gaps suggest the baseline NN is not
heavily reliant on the extra 58 dimensions --- most are redundant
one-hot columns.

\begin{figure}[t]
\centering
\includegraphics[width=\columnwidth]{results/step4/adult_convergence.png}
\caption{Step~4 convergence and training cost. Left: validation F1
vs.\ epoch (median seed). All methods plateau at similar F1 within
$\sim$15 epochs. Right: mean epochs and wall time across 5 seeds.
PCA trains longest due to delayed early stopping; ICA and RP match
baseline.}
\label{fig:step4_convergence}
\end{figure}

%%%%%%%%%%%%%%%%%%%%%%%%%%%%%%%%%%%%%%%%%%%%%%%%%%%%%%%%%%%%%%%%%%%%%%%%%%%%%%%%
\section{Step 5: Neural Network with Cluster Features}

Cluster-derived features were extracted from the raw-data Step~1 models
(K-Means $k{=}2$, GMM $k{=}19$) in four forms: one-hot assignments,
centroid distances, GMM posteriors, and GMM one-hot. These were tested
both appended to the raw 88-d feature set and by themselves.

\subsection{Appended Features}

\begin{table}[t]
\centering
\caption{Step~5A: raw + cluster features (Adult). Bold marks the best
configuration. Continuous features (distances, posteriors) outperform
discrete one-hot encoding (OHE).}
\label{tab:step5a}
\begin{tabular}{lrccc}
\toprule
\textbf{Config} & \textbf{Dims} & \textbf{Accuracy} & \textbf{F1} & $\Delta$\textbf{F1} \\
\midrule
Baseline          & 88  & .848$\pm$.003 & .696$\pm$.004 & --- \\
+ KM OHE          & 90  & .846$\pm$.002 & .695$\pm$.003 & $-$.001 \\
\textbf{+ KM Dist}& \textbf{90} & \textbf{.845}$\pm$\textbf{.003} & \textbf{.699}$\pm$\textbf{.003} & \textbf{+.003} \\
+ GMM OHE         & 107 & .845$\pm$.003 & .697$\pm$.002 & +.001 \\
+ GMM Post        & 107 & .845$\pm$.003 & .698$\pm$.002 & +.002 \\
\bottomrule
\end{tabular}
\end{table}

K-Means centroid distances provide the largest boost (+0.003 F1,
Table~\ref{tab:step5a}). Unlike OHE, distances are continuous and
encode \emph{how far} a sample is from each centroid --- a richer
signal. GMM posteriors similarly outperform GMM OHE: soft assignment
probabilities preserve uncertainty that hard labels discard. All
appended configurations reduce seed-to-seed variance (std $\leq$ 0.003
vs.\ 0.004), suggesting that the added cluster structure gives the
optimizer a more consistent starting signal, reducing sensitivity to
weight initialization.

\subsection{Cluster Features Only}

\begin{table}[t]
\centering
\caption{Step~5B: cluster features only (Adult). GMM features alone
achieve 95\% of baseline accuracy and 82\% of baseline F1.}
\label{tab:step5b}
\begin{tabular}{lrccc}
\toprule
\textbf{Config} & \textbf{Dims} & \textbf{Accuracy} & \textbf{F1} & $\Delta$\textbf{F1} \\
\midrule
KM OHE only   &  2 & .754$\pm$.000 & .170$\pm$.000 & $-$.526 \\
KM Dist only  &  2 & .764$\pm$.008 & .370$\pm$.109 & $-$.326 \\
GMM OHE only  & 19 & .807$\pm$.001 & .567$\pm$.008 & $-$.128 \\
GMM Post only & 19 & .807$\pm$.001 & .567$\pm$.008 & $-$.128 \\
\bottomrule
\end{tabular}
\end{table}

Cluster features alone carry substantial but incomplete predictive
signal (Table~\ref{tab:step5b}). The two algorithms differ
dramatically. K-Means features (2-d) are too coarse: binary OHE
yields F1$=$0.170 ($-$0.526 vs.\ baseline 0.696), and centroid
distances improve only to F1$=$0.370 ($-$0.326). GMM features (19-d)
fare much better: F1$=$0.567 ($-$0.128), recovering 82\% of baseline
F1 and 95\% of baseline accuracy (0.807 vs.\ 0.848). This is
nontrivial given entirely unsupervised clustering that scored poorly
on geometric metrics in Step~1 (we revisit this apparent
contradiction below). GMM OHE and posteriors perform identically,
likely because with 19 components the posteriors are very peaky,
collapsing to effectively the same signal as hard labels.

\subsection{DR + Cluster OHE (Bonus)}

Fig.~\ref{fig:step5_comparison} and Table~\ref{tab:integrated}
compare all configurations end-to-end.
Adding cluster OHE to DR features partially recovers information
lost by projection (e.g., PCA + GMM: F1 from 0.687 to 0.692,
recovering $\sim$60\% of the PCA-to-baseline gap). However, no
DR+cluster combination matches the raw baseline, suggesting that the
full feature space contains discriminative detail that neither DR nor
clustering fully preserves.

\begin{figure*}[t]
\centering
\includegraphics[width=0.95\textwidth]{results/step5/adult_f1_comparison.png}
\caption{Step~5: F1 comparison across all feature configurations.
Raw + KM Distances is the best overall; cluster-only GMM features
recover $\sim$82\% of baseline F1 despite being entirely unsupervised.}
\label{fig:step5_comparison}
\end{figure*}

\begin{table}[t]
\centering
\caption{Best F1 from each step (Adult). The best overall result comes
from augmenting raw features with K-Means centroid distances.}
\label{tab:integrated}
\begin{tabular}{llcc}
\toprule
\textbf{Step} & \textbf{Best config} & \textbf{F1} & $\Delta$ \\
\midrule
Baseline & Raw (88-d) & .696 & --- \\
Step 4 & PCA (30-d) & .687 & $-$.009 \\
Step 5A & Raw + KM Dist (90-d) & \textbf{.699} & +.003 \\
Step 5B & GMM only (19-d) & .567 & $-$.128 \\
Step 5C & PCA + GMM OHE (49-d) & .692 & $-$.004 \\
\bottomrule
\end{tabular}
\end{table}

\textbf{The metrics story in retrospect.} In Step~1, K-Means
dominated GMM on every metric except NMI: higher silhouette (0.440
vs.\ 0.057), higher CH, lower DB, and even higher ARI (0.064 vs.\
0.032). Yet GMM features vastly outperform K-Means features here
(F1 0.567 vs.\ 0.170). The resolution lies in what each metric
measures. Geometric metrics (silhouette, CH, DB) reward compact,
well-separated clusters --- K-Means' two tight blobs score well, but
two clusters encode very little information. ARI measures how well partition boundaries coincide with label
boundaries. Although it is adjusted for chance, it struggles when
$k$ far exceeds the number of true classes: with 19 impure components
vs.\ a binary target, each component's weak per-class lean produces
mostly noise in pairwise comparisons, diluting the signal. NMI, by
contrast, accumulates mutual information additively across
components --- each of GMM's 19 components contributes a small amount
of label information, and these weak signals sum to a meaningful
total. At 0.096 vs.\ 0.023, NMI
correctly predicted that GMM features would carry far more downstream
classification signal than K-Means features, even though every other
metric pointed the other way.

%%%%%%%%%%%%%%%%%%%%%%%%%%%%%%%%%%%%%%%%%%%%%%%%%%%%%%%%%%%%%%%%%%%%%%%%%%%%%%%%
\section{Discussion}

\subsection{Hypothesis Resolution}

\textbf{H1 (Cluster structure reflects geometry, not labels):
Supported.} ARI $< 0.2$ everywhere (Table~\ref{tab:step1}), as
predicted. K-Means found $k{=}2$ on both datasets, corresponding to
the dominant geometric split (capital-loss vs.\ non-capital-loss on
Adult, red vs.\ white on Wine) rather than classification labels.
GMM's higher-$k$ solutions captured more label information (NMI up to
0.096 on Adult) at the cost of poor geometric separation (silhouette
$< 0.06$).

\textbf{H2 (PCA preserves the most useful signal): Largely
supported.} PCA achieved the highest NN F1 among DR methods (0.687,
$\Delta$$=$${-}$0.009 from baseline), outperforming ICA ($-$0.011)
and RP ($-$0.010), and best preserved cluster structure in Step~3
(identical $k$ and ARI to raw data). As predicted, ICA showed higher
seed-to-seed variance ($\pm$0.007 vs.\ $\pm$0.004), consistent with
its independence-optimized axes creating a less smooth loss landscape.
The surprise was RP: we predicted a meaningful F1 drop, but RP nearly
matched PCA ($-$0.010 vs.\ $-$0.009) despite 82\% reconstruction
error --- consistent with the Johnson-Lindenstrauss distance-preservation
guarantee, though stronger than expected.

\textbf{H3 (Continuous cluster features outperform discrete ones):
Partially supported.} For K-Means, the pattern is clear: centroid
distances (F1$=$0.699) outperform OHE (F1$=$0.695). For GMM, however,
posteriors and OHE perform identically (both $\sim$0.698), likely
because the 19-component posteriors are too peaky to differ
meaningfully from hard assignments. Both continuous variants (KM
Distances, GMM Posteriors) exceeded baseline F1, while KM OHE did
not; GMM OHE matched baseline within noise (+0.001). As predicted,
cluster features alone (GMM, 19-d) captured meaningful but incomplete
signal (F1$=$0.567, 82\% of baseline), demonstrating that
unsupervised decomposition recovers a nontrivial fraction of the
classification task.

\subsection{Limitations and Future Work}

The NN experiments use only the Adult dataset; Wine's multiclass
structure may respond differently to cluster augmentation. All DR
methods share PCA's $n$-components for fairness, but ICA and RP might
benefit from independent tuning. We test only two clustering algorithms with Euclidean distance;
alternative metrics (cosine, Mahalanobis) or algorithms (K-Medoids,
spectral clustering, DBSCAN) could reduce sensitivity to
zero-inflated features and reveal non-convex structure that K-Means
and GMM miss. The 5-seed design captures
optimizer variance but not data-split variance. Finally, all findings
are specific to these two datasets and may not transfer to
higher-dimensional or more imbalanced settings.

%%%%%%%%%%%%%%%%%%%%%%%%%%%%%%%%%%%%%%%%%%%%%%%%%%%%%%%%%%%%%%%%%%%%%%%%%%%%%%%%
\section{AI Use Statement}
I used Cursor (Claude) to generate and refactor experiment code, produce
plots, and assist with report drafting and outlining. I thoroughly reviewed,
verified, and understand all AI-assisted content, and all final analysis and
conclusions are my own.

%%%%%%%%%%%%%%%%%%%%%%%%%%%%%%%%%%%%%%%%%%%%%%%%%%%%%%%%%%%%%%%%%%%%%%%%%%%%%%%%
\begin{thebibliography}{99}

\bibitem{kohavi1996}
R.~Kohavi,
``Scaling up the accuracy of naive-Bayes classifiers: A decision-tree
hybrid,''
in \emph{Proc.\ 2nd Int.\ Conf.\ Knowledge Discovery and Data Mining
(KDD-96)}, Portland, OR, USA, 1996, pp.~202--207.

\bibitem{cortez2009}
P.~Cortez, A.~Cerdeira, F.~Almeida, T.~Matos, and J.~Reis,
``Modeling wine preferences by data mining from physicochemical
properties,''
\emph{Decision Support Systems}, vol.~47, no.~4, pp.~547--553, 2009.

\bibitem{jolliffe2002}
I.~T.~Jolliffe,
\emph{Principal Component Analysis}, 2nd ed.
New York, NY, USA: Springer, 2002.

\bibitem{reynolds2009}
D.~A.~Reynolds,
``Gaussian mixture models,''
in \emph{Encyclopedia of Biometrics}, S.~Z.~Li and A.~Jain, Eds.
Boston, MA, USA: Springer, 2009, pp.~659--663.

\bibitem{dasgupta2003}
S.~Dasgupta and A.~Gupta,
``An elementary proof of a theorem of Johnson and Lindenstrauss,''
\emph{Random Structures \& Algorithms}, vol.~22, no.~1,
pp.~60--65, 2003.

\bibitem{hyvarinen2000}
A.~Hyv{\"a}rinen and E.~Oja,
``Independent component analysis: algorithms and applications,''
\emph{Neural Networks}, vol.~13, no.~4--5, pp.~411--430, 2000.

\end{thebibliography}

\end{document}
